% TOP_PROBLEM: {"id":30007677,"problem_number":"LOCAL-30007677","title":"Serial acquisitions","category_id":21,"category":{"id":21,"name":"economics","display_name":"Economics","description":"Open economic theory statements, published research agendas, and proposed scoped specifications; question provenance and historical priority are qualified per record.","slug":"economics","order_index":21,"created_at":"2026-10-08T00:00:00Z"},"status":"provenance_pending","external_url":null,"legacy_ids":[],"published":false,"statement":"Estimate five-year cumulative price effects of below-threshold acquisition sequences in outpatient-provider markets, allowing nonadditivity across deals.","economics":{"original_id":166,"field":"Industrial organization and competition","kind":"empirical","formalization_basis":"proposed_specification","source_status":"not_verified","priority_status":"not_established","priority_note":"No exact open-problem passage was verified. The supporting publication does not establish priority or unresolved status.","earliest_verified_reference":null,"current_reference":{"authors":["Carl Shapiro","Ali Yurukoglu"],"title":"Trends in Competition in the United States: What Does the Evidence Show?","year":2024,"url":"https://web.stanford.edu/~ayurukog/trendsincompetition.pdf","doi":"10.3386/w32762","locator":"Section 6, printed pp. 34–36; author draft dated 24 August 2024","evidence_summary":"Calls for market-specific conduct estimates and comparisons of markup methods; more merger-efficiency and nonprice retrospectives."},"formalization_note":"The cited primary work supports the subject or supplies solutions in particular settings, but no exact open-question passage for this specification was verified. The mathematical target and restrictions are proposed here, with no canonical-conjecture or priority claim."},"statusQualification":"Provenance pending: an explicit published statement of this exact open question was not verified. This is a catalogue-authored candidate specification, not a certified open conjecture. The cited primary work supports the subject or supplies solutions in particular settings, but no exact open-question passage for this specification was verified. The mathematical target and restrictions are proposed here, with no canonical-conjecture or priority claim.","statusReviewedAt":"2026-10-08"}
% FIRST_POSTED: 2026-10-08
% ENTRY_KIND: statement_only
% !TeX program = lualatex
\documentclass[11pt]{article}
\usepackage{fontspec}
\usepackage[a4paper,margin=25mm]{geometry}
\usepackage{amsmath,amssymb,mathtools}
\usepackage{xurl}
\usepackage[hidelinks]{hyperref}
\newcommand{\sref}[2]{\hyperlink{source:#1:#2}{[S#2]}}
\setlength{\emergencystretch}{3em}
\begin{document}
% problemId:30007677
\section*{Serial acquisitions}\label{30007677}
\subsection{Definitions and mathematical statement}
Fix local United States outpatient-provider markets and a ten-year observation window. For each acquirer \(m\), let \(A_m=(a_{m1},\ldots,a_{mK_m})\) be its observed sequence of acquisitions, where \(K_m\) is the number of deals and every deal is individually below the contemporaneous notification threshold. The threshold is a descriptive selection rule, not an assumption of competitive harmlessness. Define \(P_m(A_m,h)\) as the equilibrium price index at horizon \(h=5\) years under that acquisition sequence, permitting rival entry, exit and investment. Define \(P_m(0,h)\) under the counterfactual sequence containing no acquisitions by that acquirer, with other institutional changes specified identically. The target cumulative effect is \(\tau(h)=E[P_m(A_m,h)-P_m(0,h)]\), averaged over this below-threshold acquirer population. Compare it with the sum of effects attributed to the individual deals to measure nonadditivity. Construct ownership histories covering unreported transactions and address endogenous choice of targets and acquisition timing. A valid comparison requires stated counterfactual-trends or instrument assumptions, with bounds when these assumptions are uncertain. Report heterogeneity by baseline concentration and overlapping service lines. This agenda quantifies cumulative effects in one service sector; it does not infer an effect from notification status alone or claim that all small acquisitions reduce competition.

\par\textbf{Formulation and scope.} The cited primary work supports the subject or supplies solutions in particular settings, but no exact open-question passage for this specification was verified. The mathematical target and restrictions are proposed here, with no canonical-conjecture or priority claim.

\par\textbf{Open-question provenance.} An explicit published open-question statement was not verified. Sources below are supporting literature and must not be treated as a first listing.

\par\textbf{Historical priority.} No exact open-problem passage was verified. The supporting publication does not establish priority or unresolved status.

\subsection{Short English statement}
Estimate five-year cumulative price effects of below-threshold acquisition sequences in outpatient-provider markets, allowing nonadditivity across deals.

\subsection{Sources}
\begin{itemize}\raggedright
\item[\textnormal{[S1]}]\hypertarget{source:30007677:1}{} Carl Shapiro, Ali Yurukoglu. 2024. Trends in Competition in the United States: What Does the Evidence Show?. Section 6, printed pp. 34–36; author draft dated 24 August 2024.
\url{https://web.stanford.edu/~ayurukog/trendsincompetition.pdf}.
DOI: 10.3386/w32762.
{\small\textit{Supporting or current literature; see evidence qualification. Calls for market-specific conduct estimates and comparisons of markup methods; more merger-efficiency and nonprice retrospectives.}}
\end{itemize}
\end{document}
